\subsection{Behavioral simulation of stochastic attention}
\label{methods:sc_attention}

\paragraph{Inputs, references and experimental scope.}
We evaluated the unmasked single-head mapping
$Y=\mathrm{softmax}(QK^{\mathrm T}/\sqrt{d_k})V$ from supplied $Q$, $K$ and
$V$ matrices. Unless specified otherwise, each matrix contained independent
standard Gaussian entries, with $n=16$ and $d_k=d_v=32$.
The main test set contained 128 independently generated triples (seed 220001).
Additional 128-triple sets examined $d_k=d_v=16$ and 64 (seeds 220016 and
220064), halved queries, doubled queries and correlated inputs (seeds
230001--230003). Correlated inputs were constructed as
$Q=0.7K+\sqrt{0.51}E_Q$ and $V=0.7K+\sqrt{0.51}E_V$, with independent
Gaussian $K$, $E_Q$ and $E_V$. Value-offset tests reused the main triples,
adding the scalar offsets 2, 4 and 8 to every entry of $V$.
Reference scores, softmax probabilities and output products were computed in
float64. These are synthetic operator benchmarks, not activation traces from
a trained Transformer. Learned projections, head concatenation, output
projection, masking, dropout, residual connections and layer normalization
were not simulated.

\paragraph{Frozen multivariate softmax operator.}
The joint categorical HI-SMURF checkpoint from Section 3.1 was reused without
training or fine-tuning. It has 16 states per input, TT rank 8, 16 output
categories and 16,512 parameters, trained with initialization seed 11.
Static core logits were normalized offline along each outgoing core axis;
the resulting cores were used in float64 contractions. At runtime, scores
were sorted, mean-centred and encoded as
$x_j=\mathrm{clip}((z_j+4)/8,0,1)$, with the output order restored afterwards.
Every FSM started at state 7 and followed the saturating update
$s_j(t+1)=\mathrm{clip}(s_j(t)+2b_j(t)-1,0,15)$, where the input bits were
independent Bernoulli samples of $x_j$. Following 256 burn-in cycles,
state-dependent TT contractions provided categorical distributions for $L$
output events. Exactly one category was emitted at each cycle; dividing
category counts by $L$ gave non-negative attention weights with row sum
exactly one. No runtime exponential approximation or output softmax was used
in the full stochastic core. The exponential in reference computation and
in the explicitly labelled score-only control was used only for evaluation.

\paragraph{Stochastic score computation and value aggregation.}
Before stochastic encoding, each query and key vector was divided by its
smallest power-of-two scale not less than its maximum absolute entry.
Values used an analogous scale per feature across tokens. Thus input values
were not clipped. Identical scale rules were used for every configuration.
For a scaled value $u\in[-1,1]$, the bipolar baseline generated a bit with
probability $(u+1)/2$. Score products were evaluated by bipolar bit products
(equivalent to decoded XNOR operations) and signed counts summed over
features and cycles. The baseline aggregation multiplied an independent
unipolar attention bit by a bipolar value bit, summing over keys and cycles.
The improved representation stored the sign of $u$ separately and generated
a magnitude bit with probability $|u|$. Multiplication then used the AND of
the magnitude bits and the fixed product sign. Counters were decoded using
the recorded scales and, for scores, the factor $1/\sqrt{d_k}$.
The implementation uses matrix operations on integer-valued bit arrays to
accelerate this counting simulation; real-valued input products are not
substituted for the stochastic products.

\paragraph{Stream generation and value centring.}
Arithmetic streams used either independent pseudorandom draws or distinct
coordinates of a two-dimensional scrambled Sobol net. The Sobol net used
32-bit precision, a fixed scramble seed of 7001, and independently generated
32-bit digital shifts for each operand stream. Every multiplication paired
coordinate 0 for the first operand with coordinate 1 for the second. Streams
were reused across products in the matrix operations, with independent
shifts for distinct input entries. The quasi-random sources were used only
in the arithmetic stages; FSM input and categorical-emission streams
remained independent and identically distributed. Source sequences were
regenerated for each $L\in\{64,128,256,512,1024\}$. Five random-source seeds
(101, 202, 303, 404 and 505) were evaluated for each configuration.
Value centring computed the feature-wise mean over tokens, subtracted this
mean before stochastic aggregation and added it back after decoding.
The identity in Eq.~\ref{eq:attention_value_centering} applies because the
categorical attention weights sum to one. Mean estimation, rescaling,
sorting, encoding comparisons, TT contraction and the final mean addition
were evaluated numerically; their hardware cost was not measured.

\paragraph{Calibration and controls.}
Before evaluating held-out results, a gain grid of
$\{1,1.025,1.05,1.075,1.10,1.125,1.15,1.20,1.25\}$ was tested on 128
separate Gaussian triples (seed 210001), minimizing analytic attention-output
RMSE. The selected gain was 1, so no gain correction was applied.
Development sets with seeds 210002 and 210003 were used to check the arithmetic
changes and value centring; neither set was reused in the final figures.
The primary ablation successively changed bipolar IID arithmetic to
sign--magnitude IID arithmetic, sign--magnitude Sobol arithmetic, and the
addition of value centring. Stage-isolation controls replaced only score
computation, only finite-stream HI-SMURF or only value aggregation while
keeping the other operations exact. These controls are not an additive
variance or error decomposition. The analytic control contracted the exact
IID stationary-state probabilities and used exact surrounding matrix products.
No test-set tuning or best-seed selection was used.

\paragraph{Metrics, uncertainty and numerical checks.}
For $B$ test matrices, output RMSE and relative Frobenius error were defined as
\begin{equation}
\mathrm{RMSE}=\sqrt{\frac{\sum_{b=1}^{B}
\|\widehat{Y}_b-Y_b\|_F^2}{Bnd_v}},
\qquad
\varepsilon_F=\sqrt{\frac{\sum_{b=1}^{B}
\|\widehat{Y}_b-Y_b\|_F^2}{\sum_{b=1}^{B}\|Y_b\|_F^2}}.
\end{equation}
Cosine similarity was computed between vectorized complete output matrices
and averaged across matrices. Attention error was additionally recorded as
mean row total variation distance. Reported error bars are sample standard
deviations across five sequence seeds for a fixed checkpoint and input set;
they do not estimate training-initialization uncertainty or population
confidence intervals. The scatter figure uses seed 101, chosen in advance.
All output entries from that run are plotted, without selecting tokens or
features. Categorical normalization, endpoint arithmetic and the exact
value-shift identity were verified. The steady-state checkpoint has a maximum
row-sum deviation of $1.23\times10^{-7}$ in a separate numerical check,
whereas categorical finite-stream rows sum exactly to one.
There are $3L+256$ sequential sampling-window cycles across the three stages;
this count excludes TT evaluation latency, conversion overheads and any
implementation scheduling. These simulations establish operator-level
behavior, not a continuously stochastic or gate-level hardware implementation.
